\chapter{可信机器学习基础}
\label{chap:trust-foundations}

上一部分讨论的是：模型能否从有限训练样本中学到规律，使新样本上的损失达到要求？回答这个问题，需要说明候选模型、损失函数、训练与评价数据的分布关系，以及拿什么作为比较基准。学习保证据此说明需要多少样本、可以达到怎样的风险水平，以及结论以多大概率成立。

但模型能在新样本上保持较低的预测误差，还不等于它适合承担实际任务。预测会被用来采取行动，行动又会影响人。我们还需要检查：预测错误会造成什么损失，受影响者能否要求复核，以及系统或环境变化后，谁负责发现问题、停止自动处理并补救。

一个邮件分类器在测试集上的准确率达到97\%，是否就可以让它自动删除被判为垃圾邮件的内容？如果它只负责把可疑邮件放入一个可恢复的文件夹，少量误判可能可以接受；如果它直接删除求职通知或合同附件，同样的误判率就可能造成更大的损失。即使有限样本保证支持它在某个总体上的误判率结论，也不能单独决定应当赋予隔离还是删除权限。

\term{可信机器学习}（trustworthy machine learning）研究模型及其使用过程满足合理使用要求的条件和证据。这里的可信性是相对于任务、受影响者和使用范围而言的。一个系统可以在概率预测上可靠，却泄露训练信息；也可以抵御某种攻击，却对少数群体持续产生较高的错误率。因此，可信性需要拆解为具体主张，再分别检验，而不能由一个笼统的综合得分代替。

本部分把讨论范围限定为技术可信性：考察由模型、数据、算法和技术系统证据支持的主张，并为组织审批、法律合规和公共治理提供可复核的输入。是否部署、由谁承担责任以及什么风险可以接受，仍需要有权主体结合领域规则和受影响者作出决定；技术测试通过不能替代这些判断，本部分也不声称完整覆盖人工智能治理。

要判断邮件分类器能否自动处理邮件，可以依次追问三件事：希望它满足哪些要求，现有测试和分析是否支持这些要求，以及出现哪些变化后必须重新验证。这三个问题贯穿本章。邮件例子将把它们连接成一份完整的验收记录；后面的研究讨论再说明，其他可信性问题怎样沿着这条思路展开。

图~\ref{fig:trust-technical-scope}用交叉区域标出本书的技术范围。技术证据与其他判断可以讨论同一使用情境，却没有逻辑包含关系；技术检验通过，也不能推出某次部署已经合法、获准或得到公众认可。

\begin{figure}[htbp]
  \centering
  % 交叉议题的集合示意：面积无度量意义，不表达逻辑包含或职责继承。
\begin{tikzpicture}[figure picture,foundation picture]
 \path[use as bounding box] (0,-9) rectangle (110,78);
 \draw[draw=FigureStroke,fill=FigurePaper,line width=.75pt] (75,49) ellipse (30 and 23);
 \draw[draw=FigureStroke,fill=FigurePaper,fill opacity=.4,line width=.75pt] (34,20) ellipse (30 and 23);
 \draw[draw=FigureStroke,fill=FigurePaper,fill opacity=.4,line width=.75pt] (75,20) ellipse (30 and 23);
 \draw[draw=FigureYellow,fill=FigureYellowFill,fill opacity=.65,line width=1.1pt] (34,49) ellipse (30 and 23);
 \node[align=center] at (26,59) {技术可信性\\本书范围};
 \node at (83,59) {法律合规};
 \node at (26,8) {组织审批};
 \node at (83,8) {公众治理};
\end{tikzpicture}

  \caption{技术可信性与法律合规、组织审批、公众治理的交叉范围。边界和交叉面积仅示意共享议题，不代表集合的度量、逻辑包含或职责继承。黄色轮廓标出本书着重讨论的技术证据范围，其他主体仍须依据自身适用要求作出判断。}
  \label{fig:trust-technical-scope}
\end{figure}

\section{技术可信性的范围与三个层次}
\label{sec:trust-context}

分析风险时，先沿着“预测结果、采取的行动、造成的后果”检查整个过程。例如，垃圾邮件概率只是预测结果，隔离或删除才是行动，延误收信或丢失附件才是后果。还要说明谁承担损失，以及系统准备在哪些人群和环境中使用；否则，同一个准确率数值可能对应完全不同的使用风险。

第一层是明确要求。例如，不能只说“邮件分类器很准确”，还要规定正常邮件被误拦的比例最多是多少、误拦后是否可以恢复，以及是否允许永久删除。不同要求分别涉及可靠性、鲁棒性、公平性、隐私和安全等性质。

第二层是检查证据。测试结果来自哪些邮件，测试数据与实际邮件有什么关系，理论结论依赖哪些假设，都决定了结论能够适用到哪里。压力测试可以暴露特定缺陷，机制分析可以解释缺陷的原因，人的实际使用测试则能发现操作和复核中的问题；这些证据需要结合起来看。

第三层是规定何时重新验证、出了问题如何处理。模型更新、邮件来源变化或权限扩大，都可能使原有证据不再适用。需要有人监测这些变化，并在必要时停用、回退和恢复。是否接受尚未消除的风险、是否批准使用，应由有权作出决定的人根据证据判断。

这三个层次相互依赖，却不能合并成一项指标。没有要求，测试结果不知道要支持哪种行动；没有证据，要求只是愿望；没有重新验证和失效处置的安排，系统改变后，人们仍可能错误地沿用旧测试的结论。下面先用邮件系统明确情境与危害，再分别展开要求、证据和生命周期条件。

图~\ref{fig:trust-three-layers}把三层问题与准备、复核和处置证据的人对应起来。分工不能由算法替代；实际项目还需在授权记录中写明每项责任由谁承担。

\begin{figure*}[htbp]
  \centering
  % 按用户真矢量与文楷Regular要求重建；三层职责不是算法流程。
\begin{tikzpicture}[figure picture,foundation picture,
 font=\figurefont\mdseries\fontsize{9}{11}\selectfont,
 trust row/.style={draw=FigureStroke,line width=1pt,rounded corners=1.5mm},
 trust label/.style={align=left,anchor=west,inner sep=0pt},
 trust icon/.style={draw=FigureStroke,line width=.85pt,line cap=round,line join=round}]
 \path[use as bounding box] (0,-2) rectangle (169,64);
 \node at (20,59) {层次};
 \node at (77,59) {关键问题};
 \node at (141.5,59) {责任主体};
 \foreach \y/\fillcol/\stage in {44/FigureBlueFill/属性与规范,26/FigureRedFill/证据能力,8/FigureYellowFill/失效与处置}{
  \draw[trust row,fill=\fillcol] (0,{\y-7}) rectangle (40,{\y+7});
  \draw[trust row,fill=\fillcol] (49,{\y-7}) rectangle (105,{\y+7});
  \draw[trust row,fill=\fillcol] (114,{\y-7}) rectangle (169,{\y+7});
  \node at (20,\y) {\stage};
  \draw[draw=FigureStroke,line width=.9pt] (40,\y)--(49,\y) (105,\y)--(114,\y);
  % 责任记录小图标，不以勾号暗示已经批准。
  \draw[trust icon,fill=white] (118,{\y-3.4})--(123.6,{\y-3.4})--(123.6,{\y+1.6})--(121.8,{\y+3.4})--(118,{\y+3.4})--cycle;
  \draw[trust icon] (121.8,{\y+3.4})--(121.8,{\y+1.6})--(123.6,{\y+1.6});
  \draw[trust icon] (119,{\y+.4})--(122.3,{\y+.4}) (119,{\y-1})--(122.3,{\y-1}) (119,{\y-2.3})--(121.2,{\y-2.3});
 }
 % 规范文件、证据核查、生命周期时钟，对应三个不同问题。
 \draw[trust icon,fill=white] (53,40.6)--(58.6,40.6)--(58.6,45.6)--(56.8,47.4)--(53,47.4)--cycle;
 \draw[trust icon] (56.8,47.4)--(56.8,45.6)--(58.6,45.6);
 \draw[trust icon] (54,44.4)--(57.3,44.4) (54,43)--(57.3,43) (54,41.7)--(56.2,41.7);
 \draw[trust icon] (55.5,27) circle (2.8);
 \draw[trust icon] (57.5,25)--(60,22.5);
 \draw[trust icon] (56.1,8) circle (3.2);
 \draw[trust icon] (56.1,10.2)--(56.1,8)--(58,8);
 \node[trust label] at (63,44) {必须满足什么？\\对谁、什么范围？};
 \node[trust label] at (63,26) {支持多强结论？\\还缺什么证据？};
 \node[trust label] at (63,8) {谁使主张失效？\\何时限制、回退？};
 \node[trust label] at (128,44) {任务责任方\\受影响者、领域规则};
 \node[trust label] at (128,26) {技术人员：论证\\复核者：边界};
 \node[trust label] at (128,8) {有权主体：决策\\运行责任方：处置};
\end{tikzpicture}%

  \caption{可信主张的三个层次。属性与规范明确必须满足什么，证据能力限定可以支持多强的结论，失效、处置与生命周期条件规定结论何时暂停沿用以及如何响应。三行分别对应任务责任方设定要求、技术人员论证与复核者检查边界、有权主体决定剩余风险与运行责任方处置。横线只表示同层对应；责任分工为教学示意，不能替代实际授权与审批程序。}
  \label{fig:trust-three-layers}
\end{figure*}

\subsection{从预测结果到实际后果}

模型输出通常只是行动的一个输入。设输入为$\bm{x}$，模型给出的结果为$f(\bm{x})$，使用者或系统根据规则$d$选择行动$a=d(f(\bm{x}),\bm{x})$。例如，分类器输出垃圾邮件概率，而决策规则可以选择放行、隔离或转交人工。即使$f$不变，改变$d$也会改变系统造成的后果。

为了比较这些后果，记真实状态为$y$，行动损失为$L(a,y)$。相对于部署分布$\mathcal P$，系统的期望损失为
\begin{equation}
 R_{\mathrm{sys}}(f,d)
 =\mathbb E_{(\bm{x},y)\sim\mathcal P}
 L\bigl(d(f(\bm{x}),\bm{x}),y\bigr).
 \label{eq:trust-system-risk}
\end{equation}
这个量同时依赖预测、决策规则和数据分布。用零一损失评价分类器，只计算是否预测错误；用行动损失评价系统，还要区分错误造成的代价。两种评价都可以有用，但回答的问题不同。

图~\ref{fig:trust-decision-consequences}把预测与处置分开：改变决策规则时，模型可以保持不变，行动及其后果却会改变。

\begin{figure}[!htb]
  \centering
  % 依据已确认模型设计重建的真实矢量：同一预测、两项处置、不同后果。
\begin{tikzpicture}[figure picture,foundation picture,
 font=\figurefont\mdseries\fontsize{9}{11}\selectfont,
 trust module/.style={draw=FigureStroke,line width=1pt,rounded corners=1.5mm,fill=white},
 trust wire/.style={draw=FigureStroke,line width=.65pt},
 trust icon/.style={draw=FigureStroke,line width=.8pt,line cap=round,line join=round}]
 % 简化对象图标的真实路径；调用方定义trust icon和trust wire。
\providecommand{\TrustVectorMail}[3]{%
 \begin{scope}[shift={(#1,#2)},scale=#3]
  \draw[trust icon,fill=white,rounded corners=.4mm] (0,0) rectangle (8,5.5);
  \draw[trust icon] (0,5.5)--(4,2.3)--(8,5.5) (0,0)--(2.7,2.4) (8,0)--(5.3,2.4);
 \end{scope}}
\providecommand{\TrustVectorFolder}[3]{%
 \begin{scope}[shift={(#1,#2)},scale=#3]
  \draw[trust icon,fill=FigureBlueFill,rounded corners=.4mm] (0,0)--(0,5)--(3,5)--(4,4)--(8,4)--(8,0)--cycle;
 \end{scope}}
\providecommand{\TrustVectorRecord}[3]{%
 \begin{scope}[shift={(#1,#2)},scale=#3]
  \draw[trust icon,fill=white] (0,0)--(6,0)--(6,6)--(4,8)--(0,8)--cycle;
  \draw[trust icon] (4,8)--(4,6)--(6,6);
  \draw[trust icon] (1,4.8)--(4.8,4.8) (1,3.1)--(4.8,3.1) (1,1.4)--(3.8,1.4);
 \end{scope}}
\providecommand{\TrustVectorDocs}[3]{%
 \begin{scope}[shift={(#1,#2)},scale=#3]
  \foreach \dx/\dy in {0/2,1.6/1,3.2/0}{
   \draw[trust icon,fill=white,rounded corners=.5mm] (\dx,\dy) rectangle ({\dx+6},{\dy+8});
  }
  \draw[trust icon] (4.2,6)--(8.2,6) (4.2,4.1)--(8.2,4.1) (4.2,2.2)--(7.2,2.2);
 \end{scope}}
\providecommand{\TrustVectorTrash}[3]{%
 \begin{scope}[shift={(#1,#2)},scale=#3]
  \draw[trust icon,fill=FigureRedFill] (.7,0)--(0,6)--(6,6)--(5.3,0)--cycle;
  \draw[trust icon,fill=white,rounded corners=.3mm] (-.4,6) rectangle (6.4,7);
  \draw[trust icon] (2,7)--(2,8)--(4,8)--(4,7) (1.8,1)--(1.4,5) (3,1)--(3,5) (4.2,1)--(4.6,5);
 \end{scope}}
\providecommand{\TrustVectorMesh}[3]{%
 \begin{scope}[shift={(#1,#2)},scale=#3]
  \foreach \xx/\yy/\id in {-5/-2/a,-5/2/b,0/-4/c,0/0/d,0/4/e,5/-2/f,5/2/g}{\coordinate (icon-\id) at (\xx,\yy);}
  \foreach \i in {a,b}{\foreach \j in {c,d,e}{\draw[trust wire] (icon-\i)--(icon-\j);}}
  \foreach \i in {c,d,e}{\foreach \j in {f,g}{\draw[trust wire] (icon-\i)--(icon-\j);}}
  % 层内同色、层间蓝红黄；连接保持灰色，不把整网染成单色。
  \foreach \id/\col in {a/FigureBlue,b/FigureBlue,c/FigureRed,d/FigureRed,e/FigureRed,f/FigureYellow,g/FigureYellow}{\draw[trust icon,fill=\col] (icon-\id) circle (.7);}
 \end{scope}}

 \path[use as bounding box] (0,0) rectangle (110,56);
 \draw[trust module,fill=white] (0,21) rectangle (18,38);
 \TrustVectorMail{5}{29}{1}
 \node at (9,24) {正常邮件$\bm x$};
 \draw[trust module] (26,14) rectangle (49,45);
 \node at (37.5,40) {分类器$f$};
 \node at (37.5,34) {$p=f(\bm x)$};
 \TrustVectorMesh{37.5}{23}{1.4}
 \draw[trust module,fill=FigureBlueFill] (62,34) rectangle (82,54);
 \node at (72,49) {$d_1$：隔离};
 \TrustVectorMail{68}{40}{.9}
 \TrustVectorFolder{67}{37}{1.1}
 \draw[trust module,fill=FigureRedFill] (62,4) rectangle (82,24);
 \node at (72,19) {$d_2$：删除};
 \TrustVectorTrash{65}{7}{.8}
 \draw[trust icon,fill=white] (74,8)--(77,7)--(79,10)--(76,11)--cycle;
 \draw[trust icon] (74,8)--(77,10)--(79,10);
 \draw[trust module] (88,34) rectangle (110,54);
 \node at (99,49) {原件可找回};
 \node[anchor=west,inner sep=0pt] at (90,38) {收件人};
 \TrustVectorMail{102}{36}{.75}
 \draw[foundation flow,line width=.8pt,-{Stealth[length=1.2mm,width=.8mm]}] (109,41) to[out=0,in=0] (106,44);
 \draw[trust module] (88,4) rectangle (110,24);
 \node at (99,19) {原件已丢失};
 \node[anchor=west,inner sep=0pt] at (90,8) {收件人};
 \draw[trust icon] (105,8) circle (2.5);
 \draw[trust icon] (105,9.8)--(105,8)--(106.5,7);
 \draw[draw=FigureRed,line width=.9pt] (107,4.5)--(109,6.5) (107,6.5)--(109,4.5);
 \draw[foundation flow] (18,29.5)--(26,29.5);
 \draw[draw=FigureStroke,line width=1.1pt] (49,29.5)--(54,29.5);
 \fill[FigureStroke] (54,29.5) circle (.8);
 \draw[foundation flow] (54,29.5) to[out=70,in=180] (62,44);
 \draw[foundation flow] (54,29.5) to[out=-70,in=180] (62,14);
 \draw[foundation flow] (82,44)--(88,44);
 \draw[foundation flow] (82,14)--(88,14);
\end{tikzpicture}%

  \caption{同一封正常邮件被误拦后的两种处置。预测$p=f(\bm x)$保持不变，上方隔离保留原件并允许找回，下方永久删除原件。末端收件人承担找回成本或错过期限的风险；真实状态$y$、发现延误与恢复成本仍须纳入$L(a,y)$。图示为构造情景，不保证隔离后的邮件会及时找回。}
  \label{fig:trust-decision-consequences}
\end{figure}

\begin{example}[准确率与行动损失]
\label{ex:trust-mail-cost}
考虑一个人为构造的1000封邮件测试集。模型甲误拦10封正常邮件、漏过20封垃圾邮件；模型乙误拦2封正常邮件、漏过40封垃圾邮件。甲的准确率为97\%，乙为95.8\%。若每次误拦损失为10，每次漏过损失为1，则甲的总损失为$10\times10+20=120$，乙为$2\times10+40=60$。在这一损失设定下，准确率较低的乙反而更适合使用。
\end{example}

示例中的代价由教学设定给出，并非邮件任务的普遍取值。实际使用还可能需要考虑延误、申诉和恢复成本。若隔离后的邮件可以及时找回，损失可能小于永久删除；若用户不知道发生了误拦，可恢复性也未必转化为实际补救。这要求评价范围覆盖完整使用流程。

当模型输出概率时，决策规则可以由损失比较推导。一般的条件风险与Bayes行动已经在
\BookSectionRef{sec:decision-bayes-rules}{第一卷“策略的价值与优化”章的“单步价值下的行动选择”一节}
建立；这里用同一方法比较邮件的处置方式。设$p$确实表示给定输入为垃圾邮件的概率，误拦正常邮件的成本为$c_{\mathrm{fp}}>0$，放过垃圾邮件的成本为$c_{\mathrm{fn}}>0$，正确处理的成本暂设为零。隔离与放行的条件期望损失分别为$(1-p)c_{\mathrm{fp}}$和$pc_{\mathrm{fn}}$，因此选择隔离的条件是
\begin{equation}
 p>\frac{c_{\mathrm{fp}}}{c_{\mathrm{fp}}+c_{\mathrm{fn}}}.
 \label{eq:trust-cost-threshold}
\end{equation}
相等时两种行动在当前损失模型下等价。阈值不是固定的0.5：误拦越昂贵，支持隔离所需的概率就越高。这里的$p$是给定输入后的真实条件概率；把模型估计值代入，只能得到依赖估计质量的决策。即使若干概率区间中的平均预测与实际频率相符，也还不能断言每个输入的概率都准确。第~\ref{chap:trust-reliability}章将区分这些可靠性要求，并说明如何检验它们。

\begin{example}[增加复核选项后的决策]
\label{ex:trust-review-cost}
在一个教学设定中，取$c_{\mathrm{fp}}=10$、$c_{\mathrm{fn}}=1$。若人工复核可以无误确认标签，每次固定成本为0.2，则三项候选损失为$10(1-p)$、$p$和0.2。$p<0.2$时放行更合算，$p>0.98$时隔离更合算，$0.2<p<0.98$时选择复核；两个边界处分别有两种行动成本相同。这个区域来自行动成本的比较，而不是模型额外输出了一个“犹豫”类别。
\end{example}

无误复核只是用于说明机制的简化条件。现实中人工也会出错，复核成本可能随队列长度变化，此时第三项必须包含这些后果。模型辅助决策的价值，应与可行替代方案在同一条件下比较，例如完全人工处理、简单规则或较低权限的自动化。只比较两个模型的准确率，会遗漏它们是否优于现有流程这一更基本的问题。

成本难以精确确定时，可以报告不同成本比下的决策变化。如果两种方案只在某个很窄的成本区间内有优势，使用者就能识别结论依赖了哪项判断；如果一个方案在相关区间中都更好，选择依据更稳固。这类敏感性分析检查的是使用假设的变化，应与第~\ref{chap:trust-robustness}章对输入和环境变化的分析区分。

预测损失与行动损失之间还可能隔着一个现实过程。若某个分数用于决定谁获得培训，接受培训本身会改变后续成绩；此时，成绩不再是完全独立于决策的固定标签。只在过去获得培训的人中测得预测准确，不能直接回答把名额分给其他人会产生什么收益。需要区分“预测现状”和“比较行动后果”，并说明后者所需的额外观测或因果假设。式~\eqref{eq:trust-system-risk}把$y$视为不随当前行动变化的状态，适合邮件真实类别这样的情形。若行动改变结果，就需要描述不同候选行动下的可能结果，再比较相应损失；不能直接把历史行动下观察到的成绩当成所有行动共用的$y$。第~\ref{chap:trust-fairness}章讨论干预的长期影响时，将再次遇到这一问题。

\subsection{风险的承担者与使用范围}

系统的购买者、操作人员和受影响者不一定是同一群体。招聘系统可能帮助用人方筛选简历，却主要由求职者承受误拒后果。评价时应同时记录谁获得收益、谁承担错误、谁能够发现和纠正问题。只有开发者方便测量的指标，未必覆盖这些不同影响。

风险还取决于使用范围。语言、设备、时间、输入质量和用户群体都可能改变模型的有效性。部署说明中的“适用于某类输入”，应能对应到可检查的条件；如果遇到范围外输入，系统需要有识别、限制使用或转交的安排。第~\ref{chap:trust-robustness}章将讨论环境变化如何使已有规律失效。

失效不只来自预测错误。系统可能预测正确，却暴露了私人信息；可能忠实完成用户指令，却执行了不被允许的动作；也可能被外部内容诱导而偏离任务。把这些后果写成具体情形，才能确定后续需要什么保护。NIST的人工智能风险管理框架同样强调结合使用情境识别、测量和管理风险，但框架中的原则仍需落实为当前系统的具体要求\scite{nist2023rmf}。

\section{系统属性与规范要求}
\label{sec:trust-requirements}

明确要求时，要把“希望避免什么损失”转成可以检查的行为。例如，希望避免错过重要邮件，就需要检查哪些邮件会被误拦、用户能否发现误拦，以及恢复是否及时。随后还要检查有没有遗漏重要情境，以及不同要求是否冲突。只说“避免损失”，我们不知道怎样验收；只追求一个分数，也可能没有改善实际使用。

\subsection{可信性要求的具体表达}

“模型很可靠”没有说明可靠的对象。更明确的主张可以是：对规定分布中的邮件，自动隔离部分的正常邮件比例不超过某个阈值；或者，对指定人群给出的预测区间具有至少某个覆盖率。每条主张都需要标明评价单位、总体范围、错误事件和容许水平。

评价单位尤其容易被忽略。以单次请求为单位的失败率，与以用户一周的全部请求为单位的失败率不同。一个每次调用都很少出错的工具，在被连续调用后仍可能产生较高的累计失败概率。多步任务应同时检查单步性质与最终结果，不能把前者直接当成后者。

要求也不都适合写成平均损失。有些行为需要限制发生概率，有些要求需要对不同群体分别成立，还有些属于过程约束，例如敏感数据的访问必须受到授权。可以把目标概括为在满足约束的前提下选择模型与决策规则：
\begin{equation}
 \min_{f,d}\ R_{\mathrm{sys}}(f,d),
 \qquad
 G_j(f,d)\leq b_j,\quad j=1,\ldots,m.
 \label{eq:trust-constrained-objective}
\end{equation}
这里$G_j$是第$j$项可量化要求，$b_j$是其容许水平，$m$为要求数量。无法合理量化的程序要求仍需单独保留，不能为了套用公式而虚构分数。

\subsection{可信性要求的完整性}

模型在训练和验证中表现良好，仍可能没有满足我们在实际应用中的全部预期。常见的缺口有两类：训练和验证使用的数据与实际数据不同；训练时优化的目标与实际使用目标不同。前者可能使模型学到的规律在新环境中失效，后者则可能使模型取得较高分数，却仍造成我们希望避免的损失。两类缺口也可能同时发生，需要分别检查。

先看数据不同的情况。假设训练邮件中，垃圾邮件总带有某种页脚，正常邮件都不带这种页脚。模型既可以根据正文内容判断，也可以只看页脚。只要验证和测试邮件也保持这种关系，两种模型都可能答得很好，即使测试邮件从未参与训练，也无法据此判断模型学到了哪条规律。

为了把这个例子写得更精确，令内容特征$x_1$和页脚特征$x_2$都取0或1，训练数据中总有$x_1=x_2=y$，其中$y$为邮件标签。使用$x_1$的分类器与使用$x_2$的分类器，在同分布的新样本上都会完全正确。但实际使用时，垃圾邮件可能换掉页脚，正常邮件也可能使用同样的页脚。若只有内容与标签的关系保持，两个分类器就会给出不同结果。继续增加原来源的测试邮件，可以更准确地估计原分布上的错误率，却不能检验页脚关系改变后会怎样。

这个例子说明，训练目标和测试条件可能不足以区分实际使用行为不同的模型。这类问题称为\term{欠规定性}（underspecification）：多个模型都满足已有评价要求，但评价没有约束我们真正关心的行为差异。D'Amour等展示了不同机器学习流程中的此类现象。同一来源、未参与训练的测试数据上，模型得分接近，并不意味着换到实际使用环境后，它们仍可相互替代\scite{damour2020underspecification}。

要区分上述两条规律，可以在训练中加入页脚与正文关系不同的环境，也可以专门测试更换页脚后的邮件。结构约束和领域知识还能帮助明确哪些特征应当影响判断。这些做法需要与实际任务相符。一次重新训练得到的模型通过了这样的测试，也不能保证以后用其他随机种子训练出的模型都会通过。

再看目标不同的情况。训练时如果把所有分类错误都记为一次错误，那么误拦一封普通通知与误拦一封有截止日期的合同，会受到同样的惩罚。但实际使用者可能更在意后者，因为错过合同会带来更大损失。模型即使提高了总体准确率，也未必减少了重要邮件的误拦。问题在于训练和评价只区分“正常”与“垃圾”，没有检查正常邮件内部那些错误代价不同的类型。

当这些子类型没有在标签或评价分组中体现、某些子类型上的失败又被总体指标掩盖时，就形成了\term{隐藏分层}（hidden stratification）问题。同一标签内的某些子类型可能更难识别，或者一旦识别错误就造成更严重的后果。Oakden-Rayner等的医学影像研究讨论了低频子类型、标注质量及模型依赖无关线索等问题，说明总体表现良好仍可能掩盖重要子集上的失败\scite{oakdenrayner2019hidden}。

因此，评价邮件模型时，可以根据实际需要把普通通知、有截止日期的合同等分开检查。这些按条件选出的样本组也称为数据切片；分组的理由应当是它们在任务中确有不同意义，而不只是某个聚类算法把它们分开了。发现某组错误较多后，还需检查标签质量、样本量和估计的不确定性，并用独立数据确认：反复搜索许多组后找到的最差结果，可能只是抽样波动。鲁棒性与公平性两章将进一步讨论按环境和群体分组的评价。它们共同提醒我们：总体指标没有检查到的要求，不能视为已经满足。

\subsection{不同性质之间的联系与取舍}

本部分讨论的性质对应不同问题。可靠性检查预测及其不确定性表达是否有依据；鲁棒性考察输入或环境变化后的有效性；公平性关注机会、错误与负担的分配；目标对齐与行为约束考察学习到的行为是否符合实际要求。可解释性提供理解和质疑这些行为的证据；隐私保护限制数据使用中的信息泄露；对抗安全进一步考察有人主动操纵系统时，这些边界能否保持。

这些性质可以相互支持，也可能产生取舍。例如，把更多输入转交人工可能降低自动处理部分的错误率，却减少自动处理比例；更严格的隐私保护可能增加统计估计的噪声；限制工具权限可能减少损害，也可能使任务无法完成。评价应报告同一使用目标下的变化，而不能只展示被优化的那一个指标。

将不同指标加权求和，还隐含了允许相互补偿的假设。如果未经授权披露数据属于不可接受行为，较高的准确率就不能自动抵消它。式~\eqref{eq:trust-constrained-objective}将约束与优化目标分开，正是为了明确这种区别。阈值和优先级需要由使用目的、受影响者与领域知识共同确定，学习算法只能在给定要求下求解。

\section{证据能支持什么结论}
\label{sec:trust-evidence}

要求明确后，还需区分证据能支持多强的结论、如何取得有效证据，以及如何保留复核条件。这三个问题分别对应证据的含义、产生和记录。系统改变后原证据是否仍然适用，则属于下一节讨论的生命周期条件。

假设我们要批准邮件模型自动隔离邮件，一个准确率分数还不能支持这个决定。审批者需要知道系统会做什么、可能伤害谁、证据是否支持规定要求，以及出错后怎样恢复。把这些理由和证据组织成可复核的记录，就是本章所说的\term{可信保证}（trustworthiness assurance）。这里的“保证”是一份有条件的论证，不是承诺永远不会出错。记录可以分为十项：情境与危害、要求、显式假设、论证、多源证据、反证与剩余不确定性、责任主体、风险接受、失效触发和处置。下面按它们在判断中的作用展开。

情境与危害说明任务、受影响对象和从预测到行动的过程。要求则写清验收标准，包括比例以哪些对象为分母、容许阈值是多少，以及结论适用于什么范围。

证据部分需要继续交代显式假设、论证、多源证据，以及反证与剩余不确定性。显式假设覆盖抽样关系、标签、权限、资源和替代方案；论证说明这些条件为何能支持要求，并保留量词与比较基准；多源证据可以来自理论、随机抽样、压力测试、机制分析、人因评价和监测；反证与剩余不确定性则记录失败案例、未覆盖情境、估计误差与替代解释。

最后四项把技术结论接入实际决策与生命周期管理。责任主体说明谁准备证据、监测并执行
处置，也记录谁有权作出外部决定；风险接受记录该主体在什么范围、期限和不可补偿约束下
接受剩余风险，但风险接受本身不是技术证明的结论。失效触发指出哪些数据、模型、权限、
流程变化或监测结果会暂停沿用原结论；处置则规定降级、停用、回退、通知、补救、复验与
重新开放。失效来源分别进入情境与危害、显式假设和失效触发；存在主动攻击者时，还要
写明其知识、权限、资源、目标以及可能形成的损害路径。

\subsection{理论保证与经验检验}

理论保证说明在特定假设下可以证明什么。理解这样的保证，要看结论对哪些对象成立、涉及哪些随机过程，以及允许多大的失败概率，不能只看它是否附有证明。例如，一个覆盖保证可能对校准数据和未来样本的联合抽样成立，却不保证每个固定输入都具有相同覆盖率。使用保证时，需要检查数据生成方式、算法实现和实际评价对象是否与定理一致。

经验检验观察具体样本与攻击条件下的表现。它能够发现理论假设与现实之间的差距，也能够比较没有完整理论刻画的方法，但未观察到失效不等于证明失效不可能发生。下面的计算说明，即使评价设计很简单，也需要区分这两种结论。

\begin{example}[零次失败的含义]
\label{ex:trust-zero-failures}
设固定系统面对$n$个独立同分布试验，每次失败概率为$p$。观察到零次失败的概率是$(1-p)^n$。令$(1-p)^n=\delta$，可得到零失败结果对应的单侧置信上限
\[
 p_{\mathrm{up}}=1-\delta^{1/n}.
\]
在95\%置信水平下，100次零失败对应的上限约为2.95\%，1000次零失败对应的上限约为0.30\%。这里的95\%描述构造上限的统计程序在重复抽样中的覆盖性质，不是给固定但未知的$p$赋予95\%的后验概率。这些计算要求试验单位与抽样假设成立；反复尝试直到出现一批零失败结果后再报告，不能继续沿用同一解释。
\end{example}

\subsection{检验设计与证据有效性}

\term{行为测试}（behavioral testing）将“模型应当如何响应某种输入变化”写成可以执行的检查。Ribeiro等提出的CheckList把语言能力与测试类型结合，使评价不只依赖随机测试集的总准确率。它所组织的最小功能、保持不变和定向变化等检查，帮助开发者把对行为的预期变成具体用例\scite{ribeiro2020checklist}。

沿用邮件任务，可以从这三类预期构造测试。最小功能检查给出简短且含义明确的例子，例如独立的会议通知；保持不变的检查只改动大小写或无关排版，要求正常邮件仍被放行；定向变化则在任务已将某类未经请求的推销认定为垃圾邮件的前提下，逐步加入明确的推销内容，检查垃圾邮件分数是否朝预期方向变化。每项测试都依赖任务语义。将“转发诈骗警示中的诈骗原文”一律当成诈骗，恰好会暴露模型没有区分引用与实际意图。

行为测试与随机抽样检验各有职责。前者可以有目的地暴露机制缺口，后者可以在抽样假设下估计某个总体的风险。一个压力测试集故意大量收集困难样本，其失败率一般不能直接当作部署失败率；反过来，在大规模普通样本上未发现某个低频缺陷，也不足以免除针对它的检查。只有说明样本怎样取得，数值才有相应解释。

模型选择、阈值调整和最终评价承担不同职责。如果根据测试结果反复修改模型，原来的测试集就逐渐成为开发数据。需要保留未参与选择的数据，或者使用能处理多次选择的统计方法。第~\ref{sec:reliability-guarantee}节将给出有限候选程序的一个简单保证。

比较两个系统时，应尽量保持任务、样本、资源预算和评价口径一致。对于随机生成系统，还需要记录生成次数、解码设置和选择规则。只挑选最好的回答与另一个系统的第一次回答比较，会把额外计算和选择机会混入所谓模型优势。

人工或模型评判者也属于测量过程。需要检查其判断标准、重复一致性及与实际任务后果的关系。一个偏爱长回答的评判者，可能提高冗长答案的得分，却没有提高事实正确率。评价证据应保存模型版本、数据来源、评价规则和必要的运行条件，使其他人能够复核结论。

\subsection{证据的记录与复核}

一个测试结果之所以能够被正确使用，往往依赖没有写在分数中的信息。测试数据是否包含重复记录、标签来自谁、哪些样本被排除、与预训练语料是否重合，都会影响结论。Gebru等提出的Datasheets for Datasets主张随数据提供动机、组成、收集过程与推荐用途等说明，使数据使用者能够判断它与当前任务是否相容。这是一种记录数据条件的方式，本身并不证明数据没有偏差\scite{gebru2021datasheets}。

模型也需要对应的说明。Mitchell等提出Model Cards，强调记录预期用途、评价过程以及相关群体或条件下的表现。它把一个笼统的模型名称具体化为具有版本、范围与证据的对象。模型卡可以承载本章所说的可信主张，但如果缺少真实评价，填写完整的卡片仍不能代替证据\scite{mitchell2019modelcards}。

二者的关系可以通过一次复核理解。邮件模型的记录称其正常邮件误拦率较低，但数据说明显示测试样本全部来自某一种语言和某个用户群体。复核者据此能发现，结论尚未覆盖另一种语言，而不是误把模型卡中的一个平均数理解为普遍承诺。如果之后补充了新的测试，更新的应当是具体覆盖范围和证据，不能只把模型版本号改掉。

复核还要分清想确认什么。重新运行同一程序，可以检查计算过程和随机波动；换一批独立样本，可以检查结论是否只依赖原来的样本；换到新的环境，则是在检查原结论能否用于那里。这三种检查的范围逐步扩大，但通过前一种并不保证通过后一种。公开代码有助于重现计算，对新样本和新环境仍需取得适合的数据。

对于生成系统，保存一个最终回答尤其不够。模型版本、系统提示、检索内容、工具权限、采样温度、重试次数与答案选择方式都可能改变输出分布。若比较时名义上只考察模型变化，实际却同时增加了重试机会，观测到的收益应归于整个过程，不能全部解释为模型能力提高。这些条件应随评价结果一起保存，也应出现在之后的重新验证中。

\section{失效、处置与生命周期条件}
\label{sec:trust-lifecycle}

前两层说明“希望系统做到什么”和“现有证据支持到哪里”，还没有回答这些结论怎样随系统一起使用。生命周期条件需要记录哪些参与者、环境变化或系统组件可能使结论失效，谁有权接受剩余风险，哪些观测触发限制使用，以及怎样恢复到可以再次评价的状态。本章说明审批、监测和复查需要哪些技术材料，但不替组织规定法定职责，也不替受影响者决定什么后果可以接受。

\subsection{失效来源、责任主体与风险接受}

数据来源变化、标注流程失真、供应商更新、权限扩大、人员负荷和组件故障，都可能破坏
保证所依赖的条件，本章把它们统称为失效来源。需要记录变化能够作用在哪个对象、通过
什么路径影响要求，以及现有监测能否发现。只有在有人主动选择这些变化以造成损害时，
才进一步建立威胁模型，并加入攻击者的知识、权限、资源和目标。

\term{责任主体}（accountable role）是对保证中的某项决定或行动负有明确职责的角色。准备评价、独立复核、批准使用、运行监测和事故处置可以由不同角色承担；只写“项目组负责”会让触发后无人知道谁能停用系统。\term{风险接受}（risk acceptance）则是在已知证据边界和剩余不确定性的前提下，由有权主体决定某项残余风险是否能在规定范围与期限内承担。它不是证明风险消失，也不能让一个指标的收益抵消已经声明为不可接受的约束违反。

\subsection{变更触发与组合后的重新验证}

一个组件满足要求，并不意味着组件组合后仍满足同样要求。检索器返回的内容可能改变生成模型的输入分布，自动重试可能扩大隐私泄露机会，人工审核也可能因告警过多而失效。重新验证应围绕这些新关系展开，而不只是重复组件测试。

可信主张因而需要附带\term{失效触发}（invalidation trigger）：哪些数据变化、权限扩大、模型更新或监测结果会暂停沿用原结论。记录这些条件并不削弱结论，反而使使用者知道何时能够依赖它。组织如何审批、发布和追责属于更完整的治理流程；技术材料至少要暴露决策点、所需证据和触发条件，使这些流程能够作出有依据的判断。

系统组合还可能改变错误之间的依赖关系。假设在同一垃圾邮件总体上，两个检测器各有10\%的漏检率，并且任一个报警就能拦截邮件。只有在这一总体上两个漏检事件独立的条件下，才能把系统的共同漏检率算成1\%。若它们使用同一表征并遗漏同一类邮件，共同漏检率可能仍为10\%。设置第二个检测器的价值因此要看它是否提供新的信息，而不能只看它在结构上是否“多了一层”。

人工复核也需要同样处理。人看到的可能是原始证据，也可能只是模型生成的摘要；摘要已经遗漏关键信息时，增加确认步骤未必纠正错误。若要主张人机协作降低了风险，应评价真实的协作过程，包括复核者能看到什么、需要处理多少告警，以及是否有时间采取行动。第~\ref{sec:alignment-oversight}节将进一步把这些要求落实到行为监督。

\subsection{失效处置与恢复}

触发失效后，处置目标是先限制损害，再保存足以定位问题的证据，随后恢复到可以重新评价的状态。根据危害类型，可以暂停自动动作、降低权限、切换到人工流程、回退模型或隔离数据；已经删除的内容、已经泄露的信息和仍可撤销的状态具有不同恢复边界，不能统一写成“回滚即可”。处置还应说明如何通知受影响者、如何纠正已有结果，以及哪些日志必须在满足隐私与访问约束的条件下保留。

重新开放不是原版本再次启动，而是对触发原因和修复范围重新建立证据。复验应包含已知失败的回归检查，也要检查修复是否改变正常任务、其他群体或相邻性质。责任主体据此重新判断剩余风险；若证据仍不足，就应继续限制使用或选择危害更小的替代流程。监测、处置与恢复因此不是部署后的附加步骤，而是保证成立条件的一部分。

\section{例子：邮件隔离的完整保证记录}
\label{sec:trust-mail-assurance}

前三节已经分别说明如何设定要求、评价证据并管理结论的生命周期。下面沿用章首邮件系统，把这些要素合并为一个完整的教学验收实例。模型、样本、阈值和结果均为人工构造，只用于展示怎样形成可判定的保证记录，不能当作真实邮件服务的经验数据。实例固定模型版本$v_1$及其概率修正程序$c_1$，适用于公司自有网关在批准后30日内接收的中文或英文文本邮件；加密正文、纯图像邮件和从网关外导入的邮件不在范围内。误拦可能延误合同与求职通知，漏放可能使用户接触诈骗内容，系统因此不获得永久删除权限。规则在$p<0.2$时放行，在$0.2\leq p\leq0.98$时转人工复核，在$p>0.98$时进入可恢复隔离。

独立验收集包含600封正常邮件和400封垃圾邮件。规则隔离362封，其中2封为正常邮件；转人工100封，其中70封为正常邮件；放行538封，其中10封为垃圾邮件。验收要求隔离量不少于300封、自动隔离比例不低于30\%，且以隔离邮件为分母的正常邮件比例，其单侧95\%二项Clopper--Pearson置信上界不超过2\%。当前点估计为$2/362\approx0.55\%$，置信上界约为1.73\%，自动隔离比例为36.2\%，三项均达到规定阈值。50次恢复演练从用户点击“恢复”计时，直到邮件同时出现在收件箱和搜索索引中；按最近秩规则取第$\lceil0.95\times50\rceil=48$个次序统计量，经验95\%分位数为3分钟，最大值为7分钟，分别低于5分钟和10分钟的限制。这两个数只描述当前50次演练，不构成未来恢复时延总体分位数的置信保证。

这项结论显式假定验收集与30日适用范围具有约定的抽样关系，标签可靠，$v_1$、$c_1$和两个阈值在验收前固定，人工队列可用，隔离区和搜索索引均支持恢复。论证由四个环节组成：行动损失确定三段规则，范围检查限制自动处理对象，非零最低隔离量与置信上界控制抽样不确定性，恢复演练检验补救路径。除总体计数外，中英文各有500封邮件，每组隔离181封并误拦1封；40项预先规定的行为压力测试也全部通过。

这些证据仍留下明确缺口。10封垃圾邮件被放行，新式模板、隐藏子类型、共同训练盲点和人工复核错误仍可能使系统失效；低频重要邮件的估计区间可能较宽，目前也没有真实上线监测结果。评价人员负责冻结验收集并计算上界，产品责任人批准范围与有效期，运行人员负责监测、停用和回退，用户则能够查看隔离结果并发起恢复。根据上述结果，教学实例中的产品责任人接受$v_1$在规定范围内运行30日，只允许可恢复隔离，不接受永久删除；隐私或安全约束失败时，也不得由较高准确率抵消。

模型、修正程序、阈值、标签规则或工具权限一旦改变，邮件超出声明范围，隔离量低于300封、比例低于30\%、置信上界或恢复时延越界，人工队列不可用，或者出现攻击与泄露迹象，原结论就停止适用。此时应暂停自动隔离并转入人工流程，回退到已验证版本，恢复误拦邮件并保留事件记录；补充评价和处置证据，经责任主体重新接受后再逐步开放。

\section{研究内容}
\label{sec:trust-research-content}

\subsection{发展历史}

可信机器学习没有单一的创立时刻。它由统计预测、决策理论、模式识别、信息安全、人机交互
和软件保障中的多条路线逐渐汇合而成。早期研究大多分别回答一个局部问题：预测误差怎样
估计，概率分数是否可信，模型遇到噪声会怎样，数据能否保密，或者人能否理解一项决定。
当学习系统开始参与筛查、推荐、医疗和资源分配，这些问题才被放到同一个使用过程中检查：
模型给出的数值怎样触发行动，行动影响谁，系统变化后原有证据是否仍然成立。

概率预测的评价为这段历史提供了一个较早入口。Brier分数把概率预报与实际结果之间的偏差
写成适当评分，使概率输出不只按最终类别是否命中来评价\scite{brier1950}。这项思想后来
与校准、风险估计和决策阈值汇合：同样的分类准确率可以对应不同的概率质量，同样的概率
质量在不同损失和权限下又可能导致不同决定。可信性问题由此从“预测是否正确”向“输出具有
什么含义，以及使用者可以据此采取什么行动”扩展。

\Needspace{16\baselineskip}
机器学习系统最初多以平均预测误差评价，研究重点是模型能否在与训练数据相近的新样本上
保持性能。随着概率输出被用于筛查、排序和资源分配，仅报告正确率逐渐显得不够：系统还要
说明概率是否与事件频率相符，不确定性是否能识别困难输入，以及预测区间或集合是否具有
明确覆盖含义。神经网络校准研究表明，分类性能提高并不自动带来可靠的概率表达；共形预测
则从可交换样本的排序关系出发，为预测集合建立有限样本覆盖保证
\scite{guo2017calibration,angelopoulos2022conformal}。可靠性研究由此从“模型答对多少”
扩展到“模型怎样表达不知道，以及这种表达能支持什么决策”。

这项变化也把评价单位从一个类别扩展到概率、集合和后续动作。概率校准比较同一置信水平下
事件发生的长期频率，共形预测在明确的可交换条件下控制预测集合覆盖；二者都不直接保证
单个答案正确。实际系统还要同时检查集合大小、拒答比例与错误代价，否则可以用始终给出
很宽集合或频繁转人工的方式满足覆盖，却没有提供足够有用的决策信息。

可靠性研究随后把“知道得不够”转成可操作的处置。模型可以给出集合、区间或拒答信号，
系统再据此请求更多信息、调用其他工具或转交人工。这个变化看似只是增加了一个输出选项，
实际上改变了评价分母：自动处理部分的错误率、全部请求的覆盖率和转交流程的代价必须一起
报告。于是，不确定性估计开始与选择性预测、序贯决策和人机协作相连，研究对象也从单次
预测扩展到包含备选流程的完整决策程序。

这条路线还暴露了理论保证与使用解释之间的距离。校准通常是对一组预测的频率性质，覆盖
保证也依赖规定的抽样或可交换条件；它们都不能直接说明一个固定输入上的答案正确。随着
模型输出进入高代价决策，研究者越来越重视保证中的总体、分母和随机来源，并把“未观察到
错误”“平均而言校准”和“每次决定都有依据”明确区分。后续可靠性研究正是在这一差别上
发展出风险控制、拒答和变化环境中的覆盖问题。

\Needspace{10\baselineskip}
同一时期，研究者开始更系统地考察训练条件之外的变化。常见损坏基准把噪声、模糊和天气
等变化组织成可重复测试，跨来源数据集则按医院、时间、地理区域或设备划分训练与评价环境
\scite{hendrycks2019corruptions,koh2021wilds}。经验压力测试用于发现规定变化下的失败，
随机平滑等方法尝试为局部扰动范围提供认证\scite{cohen2019smoothing}。鲁棒性研究因此
逐渐区分自然变化、分布偏移和主动对抗：三者都可能降低性能，但允许变化、攻击者能力和
证据强度并不相同。

受控损坏基准固定变化类型与强度，便于比较模型在同一压力下怎样退化；跨来源数据则保留
真实采集流程中的多项差异，更接近部署环境，却较难指出究竟哪个因素造成失败。局部认证
回答给定输入邻域内输出是否保持不变，不能替代对医院、时间或人群变化的检验。研究由此
形成互补证据：基准用于重复比较，分组外推用于暴露来源差异，认证用于证明规定集合内不存在
某类反例，三者各自覆盖不同的问题。

鲁棒性之所以逐渐成为独立问题，是因为同分布测试把许多变化排除在评价之外。采集设备、
人口构成、任务规则和输入质量可以在模型不变时改变风险，训练数据中的高准确率不能说明
这些关系仍然保持。研究者由此开始区分变化的来源、幅度和语义：有些变化应当保持答案，
有些变化意味着任务本身已经改变，还有些变化由对手有意构造。若不先划定这三类情形，
“对变化稳定”就可能把真正有意义的新信息也当成噪声消除。

这一历史转折还改变了测试的作用。随机留出样本用于估计约定总体上的平均风险，压力测试
主动收集困难情境，认证则试图排除规定集合内的反例。三者产生的数字不能直接互换，却可以
共同说明系统在哪里失效、证据覆盖到哪里。后来对分布外识别、领域泛化和持续适应的研究，
正是把一次性测试继续推进到变化发现、在线更新和恢复条件。

\Needspace{12\baselineskip}
2010年代的算法公平研究把总体误差进一步拆到不同人群和个体。个体公平要求任务上相似者
得到相似处理，机会均等类标准比较给定真实结果后的错误或选择机会，反事实公平则借助因果
模型追踪群体属性改变后的预测\scite{dwork2012awareness,hardt2016opportunity,
kusner2017counterfactual}。这条路线把统计估计与伦理、法律和公共政策中的规范问题连接
起来：技术可以测量差异、证明若干指标无法同时满足或分析干预后果，却不能替有权主体决定
哪些差异合理、哪些利益不可相互补偿。

这些标准使用的比较对象并不相同。群体条件率依赖群体划分和标签质量，个体公平依赖由任务
确定的相似性，反事实标准还依赖结构因果模型及允许保留的路径。一个模型满足某个统计等式，
不能推出它已经保护了个体，也不能推出长期资源分配没有扩大差距。公平研究因此逐步把数据
测量、规范选择、因果假设和决策后的动态影响分开记录。

算法公平把可信性问题从“系统总体上表现怎样”推进到“不同人实际经历了什么”。早期模型
评价常把样本交换地汇总，群体审计却要求保留身份、标签和选择过程之间的关系；一旦系统
决定谁能进入下一环节，标签本身还可能受到过去决策影响。研究因而不再只比较一个静态差距，
而是继续追踪数据怎样形成、阈值怎样分配机会、拒答和人工复核由谁承担，以及这些决定如何
改变下一轮可观测数据。

这项发展也让技术指标的规范边界变得清楚。选择何种群体、把哪种错误视为更严重，以及哪些
差异可以接受，都不能从样本中自动学习出来。技术研究能够揭示指标冲突、估计不确定性和
干预后果，却需要把价值选择交还给有权作出决定的主体。可信机器学习由此开始同时记录数值
证据和决定这些数值如何使用的责任边界。

\Needspace{12\baselineskip}
可解释性研究也从展示模型系数扩展到复杂模型的局部行为、内部表示和训练来源。LIME、
积分梯度和SHAP分别用局部代理、基准路径和联盟分配解释预测变化
\scite{ribeiro2016lime,sundararajan2017ig,lundberg2017shap}。这些方法的广泛使用又推动
研究者区分解释是否忠实、是否完整、人能否理解以及能否改善实际判断
\scite{doshivelez2017rigorous}。反事实解释和算法补救随后把问题推进到“哪些可行行动
能够改变结果”，使可解释性与人机交互、因果推断及决策支持发生更直接的联系。

方法增多以后，“解释看起来合理”不再足以充当验收标准。局部代理要报告近似区域和误差，
路径归因要说明基准与积分路径，联盟分配要说明缺失特征怎样处理；人员实验还要区分受试者
是否只是更快完成简化任务，还是确实在真实决策中发现错误。解释若要支持补救，就必须进一步
约束可行动变量、变化成本和连带后果，避免把输入表中可修改的数值误写成现实中可执行的行动。

可解释性路线的兴起与模型复杂化直接相关。线性模型、浅层树或规则系统常能从结构本身读出
部分决策逻辑，深层网络和大型集成却把表示与计算分散到大量参数中，传统的系数阅读不再
够用。后验解释方法由此迅速发展，但它们通常只近似某个局部行为、某条路径或某种特征分配，
并不自动恢复模型的完整推理。研究焦点随之从“能否画出一张解释图”转向“解释是否如实反映模型、
由谁使用，以及是否能发现和纠正真实错误”。

进一步追问实际用途，又把解释与审计、质疑和补救连接起来。开发者需要定位训练与模型中的
问题，复核者需要判断证据是否支持一次决定，受影响者则关心能否提出异议并获得可执行的
改正途径。同一种解释很难同时满足这些目的。可信性研究因此逐步要求在给出解释以前先规定
使用者、任务和判断标准，并把人因实验与因果可行性纳入评价。

\Needspace{10\baselineskip}
隐私与安全从另一侧改变了评价边界。差分隐私用相邻数据集上输出分布的变化限制单条数据
带来的额外信息，成员推断和训练数据提取则说明，只开放模型接口也可能暴露参与事实或具体
内容\scite{dwork2014privacy,shokri2017membership,carlini2021extraction}。隐私保证需要
明确相邻关系、保护单位、查询次数和攻击者已有信息；经验攻击则用于发现特定模型与接口的
实际泄露。形式保证没有覆盖的对象仍可能被推断，某次攻击失败也不能证明不存在更强攻击。

隐私研究改变了“数据已经删除就不再有影响”的直觉。训练过程可以把单条记录的影响分散到
参数，推理接口又允许外部参与者通过重复查询累积证据；分布式训练还把原始数据、梯度、
聚合结果和模型更新交给不同参与方。保护因此从数据表访问控制扩展到训练算法、通信协议和
发布接口。差分隐私提供可组合的形式边界，经验攻击揭示具体实现和参数选择下的残余泄露，
二者共同推动研究者明确保护单位、预算如何累积以及效用损失由谁承担。

\Needspace{20\baselineskip}
对抗样本研究
显示，普通测试正确不代表输入邻域中没有可搜索的失败；投毒、后门和模型提取又把攻击面
扩展到训练数据、模型供应链和查询接口\scite{szegedy2014intriguing,
goodfellow2015adversarial,gu2017badnets,tramer2016stealing}。隐私借鉴统计推断和密码
协议，安全借鉴威胁建模与鲁棒优化，两者都要求把保护对象、观察能力和攻击预算写清楚。

安全研究还把攻击发生的阶段纳入模型。推理时扰动测试输入，训练前投毒改变学习材料，
后门让特定触发条件激活隐藏行为，模型提取则利用查询结果恢复功能或参数信息。防御若只在
固定攻击脚本下有效，可能依赖梯度遮蔽或评测盲区；可信结论需要允许自适应攻击者根据防御
调整搜索，并把攻击失败与形式化不可达明确区分。

对抗研究又把“自然发生的困难输入”与“有人寻找失败路径”分开。前者可以通过抽样和压力
测试估计频率，后者会观察系统、学习防御规则并选择最有利的输入或时机。安全评价由此不能
把一组固定攻击样本当成攻击空间本身，而要记录对手知识、权限、查询预算和最终损害。随着
模型被接入检索、插件和外部工具，威胁面也从预测函数扩大到数据供应链、权限系统和执行器，
传统机器学习指标开始与软件安全中的纵深防御和事故响应汇合。

\Needspace{20\baselineskip}
模型开始根据人的反馈完成开放任务后，研究又从静态预测扩展到行为目标和长程过程。
从行为片段比较学习奖励，到指令跟随中的示范、奖励模型和策略优化，反馈覆盖、奖励外推
和监督者能力成为可以单独检验的问题\scite{christiano2017,ouyang2022instructgpt}。这一
流程把人的判断转成可优化信号，却也增加了新的误差来源：比较样本可能没有覆盖部署行为，
奖励模型可能在策略分布变化后外推失准，较弱监督者还可能无法识别更强模型产生的隐蔽错误。

这项变化使可信性研究从“预测一个既定标签”转向“怎样形成并维持行为目标”。监督者不再
只提供正确答案，还要比较回答、说明原则、检查过程或判断长程结果；模型则可能通过优化
这些代理信号，找到监督协议没有预见的行为。研究因此开始区分能力不足、目标表达不完整、
监督信息不足和约束在新情境中失效。对齐问题并不是在训练结束后加一道内容过滤，而是贯穿
目标形成、反馈收集、策略更新、运行监控和最终执行权限。

开放任务也改变了错误的时间尺度。一次回答中的局部错误可能立即被发现，多步计划中的早期
偏离却可能在若干次工具调用后才显现；模型还会依据监督规则调整后续行为。静态测试集由此
不足以覆盖交互过程，监督与监控本身也成为可能被规避的系统组件。可扩展监督、过程评价和
权限隔离等路线，都是在回应“监督者看不完、看不懂或来不及阻止”这一共同困难。

\Needspace{18\baselineskip}
与此同时，数据集说明书、模型卡和风险管理框架推动证据从一次实验延伸到用途、版本、
责任与生命周期记录\scite{gebru2021datasheets,mitchell2019modelcards,nist2023rmf}。
可信机器学习由此逐渐形成一个跨统计学、优化、因果推断、人机交互、信息安全、软件保障
和组织治理的研究领域。

文档化工作的作用不是给系统增加一个静态标签，而是保存后来复查结论所需的上下文。数据
来源、模型版本、预期用途、分组结果和已知限制如果没有随产物一起记录，环境变化后就难以
判断旧证据是否仍然适用。风险管理进一步要求明确谁能接受剩余风险、谁负责监测，以及发现
失效后怎样限制权限、恢复状态和重新验收。

生命周期视角补上了早期基准评价缺少的时间维度。模型权重、训练数据、提示、检索库、阈值
和外部接口都会改变，任何一项变化都可能破坏原测试依赖的条件。文档若只描述发布时的模型，
就无法回答更新后哪些结论仍然有效。可信性工作因而逐渐借鉴软件保障和安全工程中的版本
追踪、变更评审、事故响应与恢复演练，把“何时重新验证”纳入技术保证本身。

基础模型与智能体进一步扩大了输出和行动范围。同一个问题可以生成许多措辞不同但含义相同
的回答，模型还可能检索外部内容、调用工具并改变持久状态。研究对象因而从固定标签、单次
预测和静态模型，扩展到开放语义、完整生成程序、连续交互和多组件系统。可信性不再只是给
模型增加一个测试指标，而要追踪信息从数据进入参数、上下文、输出和行动的全过程。

基础模型还让过去相对分离的研究路线在同一系统中相遇。检索内容既会影响事实可靠性，也
可能携带提示注入；拒答可以降低错误，却可能把等待与复核负担集中到特定人群；解释能够
支持审计，也可能泄露训练信息或被模型用来迎合检查器。单项指标越来越难代表系统行为，
研究重心因而从分别提高各项分数，转向明确它们怎样组合、冲突和共同失效。

回看这段历史，可信机器学习的范围随着使用链条逐步扩大：从平均预测误差到概率与拒答，
从同分布测试到变化和攻击，从总体表现到群体、个体与补救，从静态模型到持续更新的多组件
系统。每次扩大都源于旧证据无法回答新的使用问题，也都没有使旧问题消失。准确率仍然需要
测量，但必须放回行动损失、受影响对象、威胁条件和生命周期中解释。接下来的“当前前沿”
小节集中讨论这些历史上逐步暴露、至今仍未解决的缺口。

\subsection{当前前沿}

可信机器学习研究模型及其使用流程在什么条件下可以被依赖，内容覆盖要求表达、不确定性、
环境变化、影响分配、解释与监督、信息与行动保护以及生命周期证据。图~
\ref{fig:trust-research-problems-map}由共同判断向外展开这些研究内容：内环区分要求与
影响、证据与行为、防护与治理，中环组织八个相互制约的问题，外环列出其中部分问题下的
代表课题。下面不再逐项解释每个基本问题，只讨论当前正在推进的主要方向，以及开放生成、
多组件系统和持续更新带来的共同挑战。

\begin{figure*}[htbp]
  \centering
  % 可信机器学习研究地图：中心问题、三类条件、八个研究问题与十五项具体课题。
\begingroup
\newcommand{\trustsector}[5]{%
  \path[fill=#5,draw=black,line width=.75pt]
    (#1:#4) arc[start angle=#1,end angle=#2,radius=#4]
    -- (#2:#3) arc[start angle=#2,end angle=#1,radius=#3] -- cycle;}
\newcommand{\trustlabel}[5][]{%
  \node[font=\figurefont\mdseries#1,text=FigureInk,align=center,
    rotate=#5,transform shape] at (#2:#3) {#4};}

\begin{tikzpicture}[figure picture,foundation picture]
  \path[use as bounding box] (-82,-82) rectangle (82,82);

  % 第一层：要求与影响、证据与行为、防护与治理。
  \trustsector{90}{-45}{16}{34}{FigureBlue!78!white}
  \trustsector{-45}{-180}{16}{34}{FigureYellow!82!white}
  \trustsector{-180}{-270}{16}{34}{FigureRed!78!white}

  % 第二层：八个相互制约的研究问题。
  \trustsector{90}{45}{34}{55}{FigureBlue!58!white}
  \trustsector{45}{0}{34}{55}{FigureBlue!58!white}
  \trustsector{0}{-45}{34}{55}{FigureBlue!58!white}
  \trustsector{-45}{-90}{34}{55}{FigureYellow!62!white}
  \trustsector{-90}{-135}{34}{55}{FigureYellow!62!white}
  \trustsector{-135}{-180}{34}{55}{FigureYellow!62!white}
  \trustsector{-180}{-225}{34}{55}{FigureRed!58!white}
  \trustsector{-225}{-270}{34}{55}{FigureRed!58!white}

  % 第三层：每个类别内部顺时针展开，末项在中层收束。
  \foreach \a/\b in {90/75,75/60,60/45,45/30,30/15,15/0}
    {\trustsector{\a}{\b}{55}{80}{FigureBlue!34!white}}
  \foreach \a/\b in {-45/-60,-60/-75,-75/-90,-90/-105,-105/-120,-120/-135}
    {\trustsector{\a}{\b}{55}{80}{FigureYellow!38!white}}
  \foreach \a/\b in {-180/-195,-195/-210,-210/-225}
    {\trustsector{\a}{\b}{55}{80}{FigureRed!34!white}}

  % 中心问题。
  \fill[white] (0,0) circle[radius=16];
  \draw[draw=black,line width=1pt] (0,0) circle[radius=16];
  \node[font=\figurefont\mdseries\fontsize{8.8}{10.5}\selectfont,
    text=FigureInk,align=center,text width=28mm] at (0,0)
    {\bfseries 可信机器学习\\[1mm]
     真实使用中的系统\\何时可以被依赖？};

  % 第一层标签。
  \trustlabel[\fontsize{8.2}{9.8}\selectfont]{22.5}{25}{要求与影响}{-67.5}
  \trustlabel[\fontsize{8.2}{9.8}\selectfont]{-112.5}{25}{证据与行为}{-22.5}
  \trustlabel[\fontsize{8.2}{9.8}\selectfont]{-225}{25}{防护与治理}{45}

  % 第二层标签。
  \trustlabel[\fontsize{7.3}{8.8}\selectfont]{67.5}{44.5}{要求\\可操作化}{-22.5}
  \trustlabel[\fontsize{7.3}{8.8}\selectfont]{22.5}{44.5}{未知识别\\与转交}{-67.5}
  \trustlabel[\fontsize{7.3}{8.8}\selectfont]{-22.5}{44.5}{风险、收益\\与负担分配}{67.5}
  \trustlabel[\fontsize{7.3}{8.8}\selectfont]{-67.5}{44.5}{变化、适应\\与恢复}{22.5}
  \trustlabel[\fontsize{7.3}{8.8}\selectfont]{-112.5}{44.5}{解释、审计\\与补救}{-22.5}
  \trustlabel[\fontsize{7.3}{8.8}\selectfont]{-157.5}{44.5}{目标、监督\\与权限}{-67.5}
  \trustlabel[\fontsize{7.3}{8.8}\selectfont]{-202.5}{44.5}{隐私、安全\\与行动保护}{67.5}
  \trustlabel[\fontsize{7.3}{8.8}\selectfont]{-247.5}{44.5}{联合保证\\与持续验证}{22.5}

  % 外层标签：要求与影响。
  \trustlabel[\fontsize{6.5}{7.7}\selectfont]{82.5}{67.5}{危害情景}{-7.5}
  \trustlabel[\fontsize{6.5}{7.7}\selectfont]{67.5}{67.5}{可检验主张}{-22.5}
  \trustlabel[\fontsize{6.5}{7.7}\selectfont]{52.5}{67.5}{处置阈值}{-37.5}
  \trustlabel[\fontsize{6.5}{7.7}\selectfont]{37.5}{67.5}{校准与覆盖}{-52.5}
  \trustlabel[\fontsize{6.5}{7.7}\selectfont]{22.5}{67.5}{拒答与转交}{-67.5}
  \trustlabel[\fontsize{6.5}{7.7}\selectfont]{7.5}{67.5}{主张级核验}{-82.5}

  % 外层标签：证据与行为。
  \trustlabel[\fontsize{6.5}{7.7}\selectfont]{-52.5}{67.5}{复合偏移}{37.5}
  \trustlabel[\fontsize{6.5}{7.7}\selectfont]{-67.5}{67.5}{变化监测}{22.5}
  \trustlabel[\fontsize{6.5}{7.7}\selectfont]{-82.5}{67.5}{适应与恢复}{7.5}
  \trustlabel[\fontsize{6.5}{7.7}\selectfont]{-97.5}{67.5}{归因与反事实}{-7.5}
  \trustlabel[\fontsize{6.5}{7.7}\selectfont]{-112.5}{67.5}{机制解释}{-22.5}
  \trustlabel[\fontsize{6.5}{7.7}\selectfont]{-127.5}{67.5}{补救有效性}{-37.5}

  % 外层标签：防护与治理。
  \trustlabel[\fontsize{6.5}{7.7}\selectfont]{-187.5}{67.5}{数据影响\\与撤回}{82.5}
  \trustlabel[\fontsize{6.5}{7.7}\selectfont]{-202.5}{67.5}{自适应攻击}{67.5}
  \trustlabel[\fontsize{6.5}{7.7}\selectfont]{-217.5}{67.5}{最小权限}{52.5}

  % 统一外轮廓与层级边界。
  \draw[draw=black,line width=.8pt]
    (90:80) arc[start angle=90,end angle=0,radius=80];
  \draw[draw=black,line width=.8pt]
    (-45:80) arc[start angle=-45,end angle=-135,radius=80];
  \draw[draw=black,line width=.8pt]
    (-180:80) arc[start angle=-180,end angle=-225,radius=80];
  \foreach \a in {90,0,-45,-135,-180,-225}
    {\draw[draw=black,line width=.8pt] (\a:55) -- (\a:80);}
  \draw[draw=black,line width=.6pt] (0,0) circle[radius=34];
  \draw[draw=black,line width=.6pt] (0,0) circle[radius=55];
\end{tikzpicture}%
\endgroup

  \begingroup
    \let\RaggedRight\Centering
    \caption{可信机器学习分层研究地图。}
    \label{fig:trust-research-problems-map}
  \endgroup
\end{figure*}

\Needspace{22\baselineskip}
\paragraph{开放任务的要求表达与可靠性控制}

开放生成和智能体任务使评价单位从单个标签扩展到事实主张、检索依据、工具调用、交互轨迹
和外部状态变化。当前研究正在把行为预期写成可执行测试，并把失败案例用于发现遗漏要求
\scite{ribeiro2020checklist}。主要挑战是，测试只能覆盖已经想到的情境；自动生成大量用例
也可能重复表面变化，而没有触及语义边界、权限边界或长程后果。

可靠性研究则从整段回答得分推进到主张级核验、证据归属和行动级风险。共形预测提供覆盖
语义，\term{语义熵}（semantic entropy）尝试把措辞不同但含义相近的回答合并后衡量分歧
\scite{angelopoulos2022conformal,farquhar2024semantic}。这些方法仍难处理模型稳定复述
同一错误、检索与筛选改变输出分布，以及生成器和验证器共享盲点等情形。

拒答、预测集合和人工转交把不确定性接入决策，也会改变覆盖率、服务延迟与群体负担。当前
前沿需要联合控制自动处理风险、拒答分布和备选流程能力，并研究多轮自适应选择后风险怎样
累积。核心困难不是再产生一个置信分数，而是使“不知道”对应可执行、可复核且不会把风险
简单转移给人工或特定人群的处置。

\Needspace{22\baselineskip}
\paragraph{组合变化与长期分配后果}

部署变化通常同时涉及设备、语言、人群、时间和使用策略，单因素扰动不能代表这些交互。
跨来源数据集揭示了复合偏移下的性能差异；在EEG基础模型这一具体场景中，持续适应方法
的收益也并不稳定，更新甚至可能损害原有环境\scite{koh2021wilds,lee2026neuroadapt}。
主要挑战是在标签延迟或缺失时区分
外观变化与风险变化，并判断哪些旧证据仍可沿用。

变化还会重新分配错误、等待和复核负担。机会均等类标准与反事实标准提供不同比较方式，
也依赖不同标签和因果假设\scite{hardt2016opportunity,kusner2017counterfactual}。当前
研究从静态群体差距转向长期反馈：拒绝、推荐和资源分配会改变下一轮可观察数据，干预后的
差距既受模型影响，也受行为响应和资源约束影响。

这一方向需要把鲁棒性、公平性和服务流程放在同一时间轴上评价。研究不仅要报告变化后的
平均表现，还要记录最坏时段、恢复时间、历史任务回归，以及拒答、主动标注和人工处理负担
落在谁身上。难点在于避免一种改进把风险转移到另一环境、另一群体或测量范围之外。

\Needspace{22\baselineskip}
\paragraph{可扩展监督、机制解释与审计}

开放任务把人的示范、偏好和原则转成训练信号，但奖励模型会在策略分布变化后外推，评价者
也可能共享盲点\scite{christiano2017,ouyang2022instructgpt}。当任务长到无法直接核验，
任务分解、过程反馈和弱到强监督试图扩大可用监督；挑战在于较弱监督者可能看不出更强系统
的隐蔽错误，过程记录也未必忠实反映真实计算。

\term{可监控性}（monitorability）研究进一步考察允许观察的信息能否帮助发现偏离
\scite{burns2023weaktostrong,guan2026monitorability}。监控器一旦参与优化，就可能被
规避；检查器与被检查模型若共享训练数据、架构或提示，还可能产生相关错误。当前前沿因而
关注独立监督、随机审计、保留验收任务和抗规避监控，而不只是在静态数据上提高检查器准确率。

机制解释试图从内部表示和计算路径寻找可审计证据，稀疏特征分解提供了观察复杂模型内部
状态的一种入口\scite{bricken2023monosemanticity}。主要困难是：给某个内部表示起一个易懂的名字，还不能证明模型实际依靠它完成任务。需要通过干预等方法检查它的作用，再把内部计算与工具调用、权限检查和外部状态变化联系起来。
解释只有能够帮助发现错误、定位干预点或支持实际补救，才构成可信性证据。

\Needspace{24\baselineskip}
\paragraph{隐私、安全与智能体行动}

隐私关心信息从数据流向模型、输出和其他参与方时受到什么限制。差分隐私通过相邻数据集上
输出分布的接近程度约束单个保护单位的额外影响，成员推断与训练数据提取则从模型接口寻找
实际泄露\scite{dwork2014privacy,carlini2021extraction}。当前研究需要把保护单位、预算
组合、外部知识和重复查询共同纳入分析，并追踪敏感信息是否仍存在于原始数据、参数、检索
索引、缓存或审计副本中。差分隐私的数学保证说明规定条件下的信息泄露受到怎样的限制；实际攻击测试则检查特定接口是否已经泄露。一次攻击失败不能证明不存在其他泄露方式。

检索、长期记忆和工具接口又把外部数据变成控制入口。间接提示注入可以借助被检索内容改变
系统行为
\scite{greshake2023indirect}。签名和来源检查可以约束组件身份，却不能证明内容安全；
输入过滤也不能替代执行端的类型、权限与预算限制。主要挑战是建立覆盖数据供应链、模型
接口和实际执行器的威胁模型，并允许攻击者根据防御反馈调整策略。

智能体把攻击后果从文本错误扩大到持久状态变化。动态评测开始同时考察正常任务、有用的
第三方指令与恶意操纵，以测量安全性和过度防御之间的取舍\scite{li2026agentdyn}。更有
解释力的评价会记录攻击是否只改变文字、是否形成工具调用、调用是否被执行、状态能否恢复，
并检查突破后的恢复。当前前沿由此从静态越狱集合转向自适应攻击、多组件供应链和端到端
损害评估，同时避免过度防御使正常任务失去可用性。

\Needspace{22\baselineskip}
\paragraph{联合保证与持续验证}

真实系统由数据管线、基础模型、检索器、策略、监控器、工具和人工流程共同构成。每个组件
在各自测试上通过，不代表组合后仍满足系统要求：检索器可能把敏感信息交给生成器，生成器
可能输出正确调用但使用了过大权限，两个错误率很低的组件也可能在同一困难输入上同时失败。
当前研究正在用保证案例和结构化证据图追踪局部主张在接口处怎样传递。挑战在于组件错误
可能相关，多项要求也可能冲突；若不先检查组件间的依赖，就直接相加、相乘或汇总为加权分数，可能掩盖共同失效，也可能让某项高分抵消本来不允许违反的要求。

一次验收也会随版本和环境变化而过期。数据集说明书、模型卡与风险管理框架推动研究记录
数据来源、预期用途、分组结果、已知限制和责任安排
\scite{gebru2021datasheets,mitchell2019modelcards,nist2023rmf}。但文档只有与实际版本、
监测信号和变更流程连接时，才能支持复核。模型权重、提示、检索库、阈值、工具权限或人工
流程任一改变，都可能要求重新判断哪些证据仍可沿用。

持续保证还要求为主张设置有效期和失效触发，用线上监测发现分布、错误、攻击和服务负担
变化，并通过分阶段开放、回归测试与恢复演练控制更新风险。研究难点包括标签延迟、低频
严重事件、多个版本并存以及监测器自身失效。可信性前沿的共同目标不是为系统颁发永久标签，
而是维护一组范围明确、证据可复核、变化可发现且失败可处置的技术主张。

\section{本章小结}
\label{sec:trust-foundations-summary}

章首的邮件分类器不能仅凭准确率决定是否自动删除邮件。三个层次分别回答：允许系统做什么、现有证据为何足以支持这些行动，以及出现哪些变化后需要停用和重新验证。相应记录要写清自动处理范围、错误比例的分母、不允许违反的要求，以及监测、停用和恢复由谁负责；剩余风险是否可以接受，还需由有权作出决定的人判断。前面的教学实例据此只支持有条件、可恢复的自动隔离，并不支持永久删除。其中2\%上限以被隔离邮件为分母，与以全部正常邮件为分母的误拦率不同；两者若都影响使用决定，就必须分别记录。

不同性质回答的问题并不相同。可靠性检查概率与风险承诺是否得到数据支持，并讨论何时把输入转交其他流程；鲁棒性处理规定变化下的有效性，公平性检查机会、错误与负担的分配；目标对齐与行为约束、可解释性、隐私保护和对抗安全分别补充行为边界、解释证据、信息保护与主动攻击条件。它们可以相互提供证据，却不能由某项改进替代其余要求。

这一研究方向从平均性能评价逐步扩展到不确定性、环境变化、群体与个体影响、解释与补救、
信息泄露、主动攻击和长程行为。基础模型与智能体又把评价单位从单次预测扩展到开放生成、
多次调用和完整行动链。当前方法虽然能够为其中若干环节提供界、测试或干预，却仍需解决
组合变化、监督不足、适应性攻击和证据随版本过期的问题。

技术可信性最终形成的是一组有边界、可复核且带处置条件的主张，而不是一个综合标签。技术证据可以揭示系统在哪些条件下达到要求、还留下哪些不确定性，并为审批和复查提供输入；它不能替代组织对职责的安排、法律判断或受影响者参与的公共治理决定。

\paragraph{延伸阅读}

NIST人工智能风险管理框架适合继续理解使用情境、测量和风险处置之间的关系；阅读时应把框架原则落实为具体系统中的责任主体、证据和失效触发，而不是只保留抽象分类\scite{nist2023rmf}。

若希望进一步建立可复核的技术记录，可对照数据集说明书与模型卡。前者强调数据的动机、组成和采集条件，后者强调模型的预期用途与分组评价；二者都帮助限定证据范围，却不能替代实际测试\scite{gebru2021datasheets,mitchell2019modelcards}。
